\documentclass[a4paper,12pt]{article}
\usepackage[utf8]{inputenc}
\usepackage[russian]{babel} 
\usepackage{graphicx} 
\usepackage{amsmath}
\usepackage{amssymb}
\usepackage{float}
\usepackage[colorlinks=true, linkcolor=blue]{hyperref}
\usepackage{enumitem}

\renewcommand{\floatpagefraction}{0.99}
\renewcommand{\topfraction}{0.99}      


%\title{Лабораторная работа №4.1}
\title{ ЛАБОРАТОРНАЯ РАБОТА №4.1.
Компенсационные методы измерений}


\author{Батанин Артём Александрович}
\date{\today}


\begin{document}

\maketitle

\newpage
\tableofcontents    
\newpage            

\section{Введение}
Существует класс электроизмерительных приборов, принцип действия которых основан на сравнении измеряемой (неизвестной) величины с известной, образцовой. В качестве элемента сравнения в таких приборах используются, как правило, рабочие эталоны единиц соответствующих величин: эталон напряжения, эталон сопротивления, эталон индуктивности и другие. Этим обеспечивается высокая точность и надежность измерений. Представителями данного класса приборов являются мосты и потенциометры.



\section{Цель работы}
Цель работы: изучение компенсационных методов измерения ЭДС и напряжений, регистрация зависимости ЭДС термопары при нагревании/плавлении и охлаждении/кристаллизации олова с использованием автоматизированной системы, определение чувствительности термопары и измерение температуры различных объектов.

\section{Оборудование}
Потенциометр постоянного тока; батарея питания потенциометра; нуль-индикатор, милливольтметр; термопара; нуль-термостат; печь с тиглем; источник питания печи; регистратор ЭДС термопары (АЦП).
\section{Ход эксперимента}
\subsection{Определение типа термопары}
Для определения типа термопары на её концах была создана заведомо известная разница температур
Один конец дифференциальной термопары термостатирован в ёмкости с плавящимся льдом,  второй конец опущен в кипящую воду. 
ЭДС термопары составил 6.5мВ, что соответствует ЭДС ХК термопары из таблицы №2 (\ref{fig:Table1})  при температуре от 90 до 100 °C.

\subsection{Определение температуры плавления олова}
Второй конец  дифференциальной термопары опущен в стакан с оловом, измерения напряжения производились при помощи АЦП с частотой 1 Гц. Одновременно с началом измерений был включен подогрев ёмкости с оловом. Через некоторое время после того как весь металл расплавился подогрев был выключен. 
На \autoref{fig:U_t} представлена зависимость термо-ЭДС от времени.




\section{Анализ результатов}

Для построения зависимости температуры металла от времени данные об термо-ЭДС были согласно таблице №2 (\ref{fig:Table1}) переведены в температуру. Для промежуточных значений ЭДС и значений  вне таблицы зависимость T(E) была линейно аппроксимирована. 
Для каждой выборки j был определен номер ближайшего значения термо-ЭДС i в таблице №2 (\ref{fig:Table1}).
Тогда аппроксимированное значение температуры будет равно:
\begin{equation} 
T_j \approx E_i + T'_i \cdot (Ej - Ei) 
\end{equation}

где $T'_i = (T_i - T_{i-1})/(E_i - E_{i-1} )$
Зависимость температуры от времени представлена на \autoref{fig:T_t}

На  \autoref{fig:T_t} отчётливо видно два плато, соответствующие плавлению олова и его кристаллизации.

\section{Оценка неопределённости измерений}
Неопределённость измерений вызвана следующими факторами:  дрейф температуры водно-ледяной смеси, тепловой шум в АЦП
\begin{itemize}[label={---}] 
    \item Дискретизация значений ЭДС в АПЦ;
    \item Внутреннее сопротивление АЦП и термопары;
    \item Дрейф температуры водно-ледяной смеси относительно 0°C;
    \item Tепловой шум в АЦП;
\end{itemize}
\subsection{Оценка систематической неопределённости}

\begin{itemize}[label={---}] 
    \item Дискретизация значений ЭДС в АЦП;
    Значение термо-ЭДС после оцифровки ограничены 5 знаками после запятой. Ниже приведен рассчёт неопределённости измерения ЭДС и температуры.
    \begin{equation} 
    \begin{aligned}
	\sigma_{E} = 10^{-5}/\sqrt{12}  \text{ В}\\
	\sigma_{E} \approx 2.89\cdot10^{-6} \text{ В}\\
	\sigma_{T} \approx 3.5\cdot10^{-5} \text{ °С}
    \end{aligned}
    \end{equation}
    \item Внутреннее сопротивление АЦП и термопары - $R \approx 10^7$ Ом $r \approx 10$ Ом 
    \begin{equation} 
    \begin{aligned}
	\sigma_{E} \approx R \cdot E/(R+r)  \text{ В}\\
	\sigma_{E} \approx 1.69\cdot10^{-8} \text{ В}\\
	\sigma_{T} \approx 2.0\cdot10^{-7} \text{ °С}
    \end{aligned}
    \end{equation}

    \item Дрейф температуры водно-ледяной смеси относительно 0°C: из-за низкой теплопроводности воды в стакане со смесью водя и льда мог создаться градиент температур от 0° до +24° (комнатная температура). На момент начала эксперимента градусник показал температуру смеси +8°C. Оценка неопределённости температуры: $\bar{T} = 4$ °C - ожидаемая температура  $\sigma_{T} = 4$ °C - дрейф температуры.
\end{itemize}



\subsection{Оценка статистической неопределённости}
Тепловой шум в АЦП вызван случайными блужданиями заряженных частиц. Амплитуда теплового шума - случайная величина которая описывается распределением Гаусса. для оценки теплового шума было построено распределение значений выборок АЦП возле плато (от 210°C до 212°C) \autoref{fig:rasp}.
Характеристики распределения: $\bar{T} = 211.06$ °C $\sigma = 0.38$ °C

\subsection{Общая оценка неопределённости}
$T_{пл} =  207.0$ [4.0](0.4) °C
В квадратных скобках записана систематическая наопределённость, в круглых скобках статистическая.

\newpage 

\section{Графики и результаты}

\begin{figure}[h!]
\centering
\includegraphics[width=1\linewidth]{Tabl_1_2.png}
\caption{Таблицы значений термо-ЭДС для ХА и ХК термопар}
\label{fig:Table1}
\end{figure}

\begin{figure}[H]
\centering
\includegraphics[width=1\linewidth]{U_t.png}
\caption{Изменение термо-ЭДС}
\label{fig:U_t}
\end{figure}

\begin{figure}[H]
\centering
\includegraphics[width=1\linewidth]{T_t.png}
\caption{Изменение температуры олова }
\label{fig:T_t}
\end{figure}

\begin{figure}[H]
\centering
\includegraphics[width=0.8\linewidth]{rasp1.png}
\caption{Распределение измеренных значений температуры }
\label{fig:rasp}
\end{figure}
\clearpage


\end{document}